\PassOptionsToPackage{table,xcdraw}{xcolor}
\documentclass[11pt, a4paper]{article}

\usepackage[T1]{fontenc}
\usepackage[utf8]{inputenc}
\usepackage{tgpagella}
\usepackage[spanish, es-noshorthands]{babel}
\usepackage{amsmath, amssymb}
\usepackage[most]{tcolorbox}
\usepackage[american]{circuitikz}
\usepackage{tabularx}
\usepackage[margin=2.5cm]{geometry}
\usepackage{fancyhdr}

\pagestyle{fancy}
\fancyhf{}
\setlength{\headheight}{16pt}
\lhead{\textbf{UCB Sede Tarija} | Ing. Mecatrónica}
\rhead{IMT-221 | Guía Previa Hito 2}
\renewcommand{\headrulewidth}{0.4pt}
\definecolor{ucbblue}{RGB}{0, 51, 102}

\begin{document}

\begin{tcolorbox}[colback=red!5, colframe=red!70!black, title=\textbf{Guía Previa y Lista de Materiales (BOM) — Hito 2}]
    \textbf{Par Diferencial BJT discreto para sensado de corriente mediante Shunt[cite: 17].} \\
    \textit{Importante: Esta guía fija el circuito pedagógico propuesto. Los valores son para el montaje académico. El pinout exacto debe verificarse en el datasheet de sus componentes[cite: 17].}
\end{tcolorbox}

\section{Problema a Resolver y Arquitectura}
Queremos medir una corriente de forma indirecta ($I_{sense} \approx 5\,\text{mA}$) usando una resistencia Shunt muy baja ($R_{sh} = 1\,\Omega$)[cite: 17]. Esto produce una minúscula caída de voltaje $V_{sh} = 5\,\text{mV}$[cite: 17]. El Par Diferencial amplificará esta pequeña diferencia, rechazando ruido común[cite: 17].
La arquitectura exige tres bloques: Generador/Shunt, Par Diferencial (Q1-Q2), y Espejo de Corriente de Cola (Q3-Q4)[cite: 17].

\section{Esquema Completo Propuesto}
La alimentación nominal de diseño será $+9\,\text{V}$, $0\,\text{V}$ (Tierra Central) y $-9\,\text{V}$[cite: 17].
\begin{center}
\begin{circuitikz}[scale=0.9, transform shape, thick]
    % --- Generador Shunt (Aislado Visualmente) ---
    \draw (-3,4) node[above]{$+9\,\text{V}$} to[R, l={$R_{LIM}$}, a={$1.8\,\text{k}\Omega$}] (-3,1.5);
    \draw (-3,1.5) node[circle,fill,inner sep=1.5pt](V1){} node[left]{$v_1$};
    \draw (-3,1.5) to[R, l={$R_{sh}$}, a={$1\,\Omega$}] (-3,-0.5);
    \draw (-3,-0.5) node[circle,fill,inner sep=1.5pt](V2){} node[left]{$v_2 = 0\,\text{V}$} node[ground]{};
    
    % --- Par Diferencial ---
    \draw (1,2) node[npn](Q1){}; \node[right=1cm] at (Q1.B) {Q1};
    \draw (5,2) node[npn, xscale=-1](Q2){}; \node[left=1cm] at (Q2.B) {Q2};
    
    % Conexiones de Entrada
    \draw (V1) |- (Q1.B);
    \draw (Q2.B) -- (6,2) node[ground]{} node[right]{$v_2 = 0\,\text{V}$};
    
    % Resistores de Colector
    \draw (Q1.C) to[R, l={$R_{C1}$}, a={$4.7\,\text{k}\Omega$}] (1,4.5) node[above]{$+9\,\text{V}$};
    \draw (Q2.C) to[R, l={$R_{C2}$}, a_={$4.7\,\text{k}\Omega$}] (5,4.5) node[above]{$+9\,\text{V}$};
    
    \draw (Q1.C) node[circle,fill,inner sep=1.5pt]{} node[right]{$v_{C1}$};
    \draw (Q2.C) node[circle,fill,inner sep=1.5pt]{} node[left]{$v_{C2}$};
    
    % Unión Emisores
    \draw (Q1.E) -- (1,0.5) -- (5,0.5) -- (Q2.E);
    
    % --- Espejo de Corriente ---
    % Q4 (Transistor de Cola) apuntando hacia la derecha
    \draw (3,-1.5) node[npn, xscale=-1](Q4){}; \node[left=1cm] at (Q4.B) {Q4};
    % Q3 (Referencia Diódica)
    \draw (6,-1.5) node[npn](Q3){}; \node[right=1cm] at (Q3.B) {Q3};
    
    % Conexión Cola-Emisores
    \draw (3,0.5) node[circle,fill,inner sep=1.5pt]{} -- (Q4.C);
    
    % Unión de Bases
    \draw (Q4.B) -- (Q3.B);
    % Conexión Diódica de Q3
    \draw (Q3.B) node[circle,fill,inner sep=1.5pt]{} |- (Q3.C) node[circle,fill,inner sep=1.5pt]{};
    
    % Resistor de Referencia
    \draw (Q3.C) to[R, l={$R_{REF}$}, a_={$8.2\,\text{k}\Omega$}] (9,-0.75) node[ground]{} node[above]{$0\,\text{V}$};
    
    % Conexión Emisores del Espejo a -9V
    \draw (Q4.E) -- (3,-3) node[below]{$-9\,\text{V}$};
    \draw (Q3.E) -- (6,-3) node[below]{$-9\,\text{V}$};

\end{circuitikz}
\end{center}
\section{Cálculos Pre-Laboratorio Exigidos}
Deberá incluir las siguientes derivaciones matemáticas exactas en su reporte previo:
\begin{enumerate}
    \item \textbf{Generación de Señal:} Calcule analíticamente $I_{sense} = \frac{9\,\text{V}}{1800\,\Omega + 1\,\Omega}$ y obtenga $V_{sh}$[cite: 17].
    \item \textbf{Espejo de Corriente:} Determine la corriente de referencia. Asuma que la base de Q3 está en $-9\,\text{V} + 0.7\,\text{V} = -8.3\,\text{V}$. Derive $I_{REF} = \frac{0\,\text{V} - (-8.3\,\text{V})}{8.2\,\text{k}\Omega}$ y compruebe que resulta $\approx 1.01\,\text{mA}$. Justifique por qué $I_T \approx I_{REF}$[cite: 17].
    \item \textbf{Punto de Operación del Par:} Usando $I_T \approx 1\,\text{mA}$, asuma $v_1 \approx v_2 \approx 0\,\text{V}$. Calcule $I_{CQ}$, las caídas en los resistores de colector, el potencial de emisor $V_E$, las tensiones de salida $V_C$ y el $V_{CE}$ para verificar la región activa[cite: 17].
    \item \textbf{Ganancia Diferencial Teórica ($A_d$):} Calcule $g_m = I_{CQ}/V_T$ (use $25.85\,\text{mV}$). Utilizando $A_d \approx g_m R_C$, determine el voltaje de salida diferencial esperado $v_{od} = A_d \cdot v_d$ considerando la entrada de $5\,\text{mV}$ del Shunt[cite: 17].
\end{enumerate}

\section{Lista de Materiales Oficial (BOM)}
Se exigen los siguientes componentes por grupo[cite: 17]:
\begin{itemize}
    \item \textbf{Transistores:} 6 unidades 2N3904 (Q1-Q4 más repuestos para seleccionar el mejor \textit{matching})[cite: 17].
    \item \textbf{Resistores (Película Metálica 1\%):}
    \begin{itemize}
        \item 4.7 k$\Omega$ x 4-6 unidades (Seleccionar el par más idéntico para $R_{C1}$ y $R_{C2}$)[cite: 17].
        \item 8.2 k$\Omega$ x 2 (Para $R_{REF}$)[cite: 17].
        \item 1.8 k$\Omega$ x 2 (Para $R_{LIM}$)[cite: 17].
        \item 1 $\Omega$ ($\ge 1/4\,\text{W}$) x 2 (Para Shunt)[cite: 17].
    \end{itemize}
    \item Protoboard, jumpers y acceso a fuente dual $\pm 9\,\text{V}$ (o dos baterías de 9V cuadradas)[cite: 17].
\end{itemize}
\textit{Atención: Para lograr el CMRR $> 60\,\text{dB}$, la simetría de Q1-Q2 y $R_{C1}-R_{C2}$ es crítica. Use un multímetro para medir y seleccionar los componentes más parecidos antes de montar. Verifique el orden de los pines (E-B-C) en el datasheet de su marca de transistor[cite: 17].}

\end{document}
